\section{Exact signed data and exterior entry}\label{prep:section}

Fix $d\in\mathcal U$, a branch supplied by
Corollary~\ref{test:branch-seed}, and its preparation in
Proposition~\ref{bg:packet}.  We use the real tensors $h_T^p$ and
$h_T^B$ for the deep packet and its exterior beam realization,
respectively, and write $q_T$ for the normalized component array of
$h_T^p$.  The scales are
\[
 A_T=e^{-\alpha T},\qquad \alpha=7\kappa/8,\qquad
 \lambda_T=e^{\zeta T}.
\]
The construction in this section takes place in ordinary, local
wave-map gauge with target the actual background metric $g$.
All spatial data are eventually put on the complete original bridge.
Fix an unbounded sequence of admissible late branch levels $T$,
with the fixed launch and angular-atlas margins used in
Proposition~\ref{bg:packet}, and the corresponding packets.
Limits in this section are taken along this sequence.

Put $\mathcal S_T=\mathcal S\cap\mathcal D_T$, where
$\mathcal S=\{t=0\}$ and $\mathcal D_T$ is defined in
Equation~\eqref{bg:packet-domain}. The constructions first use the
fixed slight enlargements specified there, so that all entry jets
are defined on a spacetime collar of $\mathcal S_T$.
For a spacetime tensor $W$ defined near this interval, let $\widehat W$ denote
its components in the background normalized frame of
Equation~\eqref{bg:frames}.  Define the full entry-jet norm by
\begin{equation}\label{prep:entry-norm}
 \mathfrak J_{r,T}(W)=
 \sum_{|I|\le r}\left(
   \left\|D^I\widehat W\big|_{\mathcal S_T}\right\|_{L^\infty}
  +\left\|D^I\widehat W\big|_{\mathcal S_T}\right\|_{L^2(dx\,d\operatorname{vol}_\gamma)}
                         \right).
\end{equation}
The sum uses a fixed finite angular atlas and includes ordinary
spacetime derivatives, including derivatives normal to $\mathcal S$.
It is not a norm of tangential jets alone.  We also write
$\mathfrak J_{r,T}(F)$ when $F$ is already the normalized array.
The entry interval has length $O(T)$ and uniformly controlled angular
volume, so its $L^2$ part costs at most a polynomial factor relative
to the corresponding uniform jet bound.

\begin{proposition}[Signed preparation and entry]\label{prep:signed-entry}
Fix a finite data order $m\ge10$, a finite entry order $r$, a loss
$\eta>0$, and $0<\nu_0<1$.  After fixing all the finite derivative
orders needed below, one can choose $\zeta>0$ and then the finite
formal accuracy orders in Proposition~\ref{bg:packet} so that the
following hold for all sufficiently large parameters $T$ in the
packet sequence.

There are smooth positive complete vacuum data $d_T^\pm\in\mathcal D$
on the whole bridge, with
\begin{equation}\label{prep:data-convergence}
 p_m(d_T^\pm-d)\longrightarrow0.
\end{equation}
In particular $d_T^\pm\in\mathcal U$ for large $T$.
They agree exactly with $d$ on the other end, the connecting region,
and a protected region containing every fixed compact subset for
large $T$.  On the selected asymptotic end the protected ball is
\begin{equation}\label{prep:protected-ball}
 B_T=\{|z-z_*|<\rho_T\},\qquad
 z_*=C_e(t_T),\qquad \rho_T=(1-2\beta)t_T,
\end{equation}
where $\beta>0$ is fixed sufficiently small.  The data have the
weight-one metric and weight-two second-form symbol bounds at every
order.  Their full correction tails are retained; asymptotic
coefficients are allowed to change.

Ordinary exterior vacuum evolutions $g_T^\pm$, in exact wave-map gauge
with the same target $g$, exist on the finite domains needed to reach
$\mathcal S_T$, with fixed overlap margins and strict future-outflow
artificial boundaries.  Set
\[
 H_T^\pm=\widehat{g_T^\pm-g},\qquad
 F_T^\pm=H_T^\pm\mp q_T
 \quad\hbox{on the entry collar}.
\]
Then
\begin{equation}\label{prep:entry-estimates}
 \mathfrak J_{r,T}(F_T^\pm)\le A_T^2e^{\eta T},\qquad
 \mathfrak J_{r,T}(F_T^+-F_T^-)\le A_T^3e^{\eta T}.
\end{equation}
The exact changes, packet, and errors vanish on the earlier needed
entry interval $v<(1-\nu_0)T$.  In particular, $r=61$ is permitted,
with an arbitrarily small prescribed loss.  The scaled linear deep
residual can simultaneously be at most $e^{-4\kappa T}$ through
order $60$, or at any other finite orders fixed in advance.
The same statement holds for the separately prepared exchanged branch.
\end{proposition}

The proof follows the data from the bridge to the entry cylinder.
First we correct the two constraint defects with one linear inverse.
Then we initialize exact wave-map gauge and evolve both corrected data
with the same target metric. Once the linear residuals are below the
cubic scale, the leading discrepancy at each step is quadratic in the
beam and has the same value for the two signs;
subtracting the discrepancies leaves a cubic bound. The final exterior
energy estimate must preserve both orders, so its growth rate is made
arbitrarily small by keeping the change away from the bounded-radius
region until a late time.

The exterior geometry and linear beam estimates are
those of Proposition~\ref{bg:exterior} and
Proposition~\ref{bg:packet}.  The underlying source assertions are
\cite[\CompanionRef{foundation}{pk2:escape}, \CompanionRef{foundation}{pk2:beam-proposition}, and \CompanionRef{foundation}{pk2:exterior-proposition}]{Foundation}.
The signed constraint correction and exterior comparison are proved below,
with their common linear inverse and common propagation rate made explicit.

\subsection{A common support-preserving constraint solve}

Localized constraint deformation has a substantial history
\cite{Corvino2000,CorvinoSchoen2006,ChruscielDelay2003,CarlottoSchoen2016}.
The inverse used here preserves an inner ball and allows a decaying
exterior tail, so asymptotic coefficients may change. Its construction
in the companion uses support-preserving divergence primitives of
Bogovskii type; their support and derivative gain are treated in
\cite{CostabelMcIntosh2010}. The following signed estimate additionally
requires that both corrections use the very same linear inverse.

For tensors in Cartesian components on $\mathbb R^3$, put
\begin{equation}\label{prep:weighted-spaces}
 \|f\|_{H^j_\ell}
 =\|f\|_{H^j(B_2)}+
   \sup_{R\ge1}R^\ell\|f(R\,\cdot)\|_{H^j(1<|y'|<2)},
 \qquad
 X^s=H^{s+2}_1\oplus H^{s+1}_2,\quad Y^s=H^s_4.
\end{equation}
The two components of $X^s$ are symmetric metric and second-form
tensors; the components of $Y^s$ are a scalar and a covector.
Let $\Phi$ be the constraint map in Equation~\eqref{phase:constraints}.
Its Euclidean linearization is
\begin{equation}\label{prep:flat-linearization}
 P_{\mathrm{flat}}(q,p)=
 \bigl(\partial_i\partial_jq_{ij}-\Delta q_{ii},\,
       \partial_jp_{ij}-\partial_i p_{jj}\bigr).
\end{equation}
We use the following precise property of
\cite[\CompanionRef{foundation}{pk:constraint-linear}]{Foundation}.  One linear operator
$\mathcal K:Y^s\to X^s$, bounded for every integer $s\ge3$, satisfies
\begin{equation}\label{prep:right-inverse}
 \supp z\subset\{|y'|\ge1\}\quad\Longrightarrow\quad
 P_{\mathrm{flat}}\mathcal Kz=z,\qquad
 \supp\mathcal Kz\subset\{|y'|\ge1\}.
\end{equation}
It is defined on all of $Y^s$, not just this support subspace.
Its conjugates under rotations and positive argument dilations are
smooth in operator norm $Y^s\to X^s$.  These properties concern the
flat constraints and have no Kerr parameter restriction.  The
construction in that lemma allows a decaying tail and imposes no
moment condition on the original source.

Here is the nonlinear and signed consequence in the form we need.

\begin{lemma}[Signed correction with one inverse]\label{prep:constraint-pair}
Let $s\ge3$.  Let $D_\pm=(\delta+a_\pm,b_\pm)$ be smooth positive
data on $\mathbb R^3$ with symbol orders $-1,-2$ at every order, and
suppose
\[
 \|D_\pm-(\delta,0)\|_{X^s}\le\varepsilon_s.
\]
Suppose the smooth compact sources $\phi_\pm$ are supported in
$\{|y'|\ge1\}$ and equal $\Phi(D_\pm)$ there.  For sufficiently
small $\varepsilon_s$ and $b_2:=\max_\pm\|\phi_\pm\|_{Y^s}$ there
are corrections $v_\pm$ supported in $\{|y'|\ge1\}$ such that
\[
 \Phi(D_\pm+v_\pm)=0\quad (|y'|\ge1),\qquad
 \|v_\pm\|_{X^s}\le C_s b_2.
\]
They are obtained with the same $\mathcal K$ and satisfy
\begin{equation}\label{prep:abstract-odd-correction}
 \|v_+-v_-\|_{X^s}
 \le C_s\left(\|\phi_+-\phi_-\|_{Y^s}
                  +\|D_+-D_-\|_{X^s}b_2\right).
\end{equation}
Each correction is smooth across the unit sphere and has symbol
bounds of orders $-1,-2$ at every derivative order.  Only the stated
base norms have to be small.  No vacuum equation is required for
$D_\pm$ in the artificial interior $|y'|<1$.
\end{lemma}

\begin{proof}
Write
\[
 N_D(v)=\Phi(D+v)-\Phi(D)-P_{\mathrm{flat}}v.
\]
The subtraction is essential: $P_{\mathrm{flat}}$ sends a general
$X^s$ element to quantities of weight three, and we do not assert
that the full constraint map has values in $Y^s$ on an arbitrary
$X^s$ neighborhood.  After its flat linear terms cancel, the scalar
remainder has the product types
\[
 a\,\partial^2a,\quad (\partial a)^2,\quad b^2,
\]
and the momentum remainder has types
\[
 a\,\partial b,\quad (\partial a)b,
\]
with smooth bounded inverse-metric coefficients.  Here factors may
be from the background or the correction, and at least one is a
correction factor.  The weight budgets are $1+3=4$ and $2+2=4$.
Unit-ball and rescaled-annulus Sobolev multiplication, valid for
$s\ge3$, therefore gives
\begin{align}
 \|N_D(v)-N_D(w)\|_{Y^s}
   &\le C_s(\|D-(\delta,0)\|_{X^s}+b)
                   \|v-w\|_{X^s},\label{prep:contraction-estimate}\\
 \|N_{D_+}(v)-N_{D_-}(v)\|_{Y^s}
   &\le C_s\|D_+-D_-\|_{X^s}\|v\|_{X^s},
                    \label{prep:mixed-constraint}
\end{align}
when $\|v\|_{X^s},\|w\|_{X^s}\le b$.  For the second inequality,
integrate the mixed second differential of $\Phi$ along the segments
in the datum and correction variables.  The expression vanishes at
$v=0$, and the same displayed weight budgets bound this differential.
Differentiation of the inverse-metric coefficients preserves these
estimates on a uniformly positive small ball.

For both signs solve
\begin{equation}\label{prep:shared-fixed-point}
 z_\pm=-\phi_\pm-N_{D_\pm}(\mathcal Kz_\pm),\qquad
 v_\pm=\mathcal Kz_\pm.
\end{equation}
On the common ball $\|z\|_{Y^s}\le2b_2$, the first inequality has
Lipschitz constant at most $1/2$ after shrinking the two base sizes.
The map preserves this ball because $N_D(0)=0$.  Starting at zero
preserves exterior support by Equation~\eqref{prep:right-inverse}
and locality of $N_D$.  The limit consequently has that support,
and $P_{\mathrm{flat}}v_\pm=z_\pm$ proves the claimed exterior
constraint equation.  Subtract the two fixed-point equations, apply
Equation~\eqref{prep:mixed-constraint} at $v_-$, and absorb the
$\frac12\|z_+-z_-\|_{Y^s}$ term.  Boundedness of $\mathcal K$
gives Equation~\eqref{prep:abstract-odd-correction}.

We include the all-order argument from the proof of
\cite[\CompanionRef{foundation}{pk:constraint-nonlinear}]{Foundation} to specify that it regularizes these
same fixed points.  Let $U_\vartheta,V_\vartheta$ be the actions on
$X^s,Y^s$ of a rotation and a positive argument dilation near the
identity; rotations also act on tensor components.  For either
sign put
\[
 \mathcal K_\vartheta=U_\vartheta\mathcal K V_\vartheta^{-1},
 \quad \phi_\vartheta=V_\vartheta\phi,
 \quad N_\vartheta(w)=V_\vartheta N_D(U_\vartheta^{-1}w).
\]
The first family is norm-smooth by the cited inverse lemma.
The all-order symbol bounds of $D$ and smooth compact support of
$\phi$ make the other displayed expressions smooth in their
parameters.  In the explicit weighted products for $N_\vartheta$,
parameter differentiation differentiates background coefficients or
the explicit component and scale factors; it does not require more
unknown derivatives than the two metric and one second-form
derivatives allowed by $X^s$.

On the full space $X^s$ consider
\begin{equation}\label{prep:regularity-equation}
 w+\mathcal K_\vartheta\phi_\vartheta
                 +\mathcal K_\vartheta N_\vartheta(w)=0.
\end{equation}
Its derivative in $w$ has the form $\Id+\mathcal K_\vartheta
DN_\vartheta(w)$ and is invertible by a Neumann series on a common
small ball.  The implicit-function theorem gives a smooth local
solution.  The already constructed $U_\vartheta v$ satisfies the
same equation algebraically.  Choose $b_2$ small enough that $v$
lies in a smaller ball; uniform boundedness of $U_\vartheta$ places
all these transforms in the common uniqueness ball.  Thus uniqueness
identifies their orbit with the smooth implicit-function branch.
This argument does not assume norm continuity of the unknown orbit.
It uses the all-input definition of $\mathcal K$, because dilation
moves the support boundary.

Distributional differentiation of the orbit now gives every iterate
of rotations $\Omega_{ij}=y'_i\partial_j-y'_j\partial_i$ and the
dilation $E=y'\cdot\partial$ in $X^s$.  Constant component rotations
contribute only removable zeroth-order matrices.  Since
\[
 |y'|^2\partial_i=y'_iE+\sum_j y'_j\Omega_{ji},
\]
these derivatives span the scaled coordinate derivatives on every
exterior annulus and on a neighborhood of $|y'|=1$.  Scaled Sobolev
embedding proves smoothness there and all differentiated symbol
bounds.  Inside the sphere the correction is zero.  Hence it is
smooth across the sphere as well.  Higher bounds may depend on the
finite experiment; their smallness has never been required.
Finally, the small base norm preserves positive definiteness.
\end{proof}

\subsection{From beam seeds to whole-bridge data}

Use the regular exterior coordinates of Proposition~\ref{bg:exterior},
with $z=y+C_e(t_e)$, $I(t_T)=T$, and
$0<c_I\le I'\le C_I$.  Thus $t_T\asymp T$.  The background drift
$s_*:=\sup|C_e'|$ is small; choose the fixed localization parameter
$\beta$ so that $1-2\beta-s_*>0$.  The beam return-shell estimate,
with a slightly smaller tolerance if needed, places the induced
beam seed strictly outside the ball in
Equation~\eqref{prep:protected-ball}.  All beam tubes have this
margin.  At $t_e=0$ this far portion of the partial exterior slice
is the original bridge.

For $-1\le\tau\le1$, induce metric and second form from
$g+\tau h_T^B$.  Away from the beam support retain $d$.  On the
selected far end introduce
\begin{equation}\label{prep:rescaling}
 y'=(z-z_*)/\rho_T,\qquad
 \widetilde h(y')=h(z_*+\rho_Ty'),\qquad
 \widetilde K(y')=\rho_T K(z_*+\rho_Ty').
\end{equation}
These are rescaled component functions, with the second expression
for the second form including the displayed factor.  Directly in
the constraint formulas, both components of $\Phi$ acquire the
factor $\rho_T^2$.  Choose fixed radii
$s_* /(1-2\beta)<r_0<r_1<1$.  Since
$|z_*|\le s_*t_T+O(1)$, the triangle inequality gives
$|z|\asymp\rho_T|y'|$ for $|y'|\ge r_0$; a bounded initial offset
of $C_e$ is harmless for large $T$.  The symbol assumptions on $d$
give, at every fixed order $j$,
\[
 |\partial_{y'}^j(\widetilde h-\delta)|
 \le C_j\rho_T^{-1}|y'|^{-1-j},\qquad
 |\partial_{y'}^j\widetilde K|
 \le C_j\rho_T^{-1}|y'|^{-2-j}.
\]
Cut these background deviations off to zero between radii $r_0$
and $r_1$.  Leave them unchanged for $|y'|\ge r_1$, and add the
scaled beam seed, whose support is a fixed annulus strictly outside
the unit sphere.  Denote the resulting artificial data by $D(\tau)$
and set $D_\pm=D(\pm1)$.

The artificial data need not be vacuum inside the unit ball.
Outside it, $D(0)$ is the original vacuum datum.  Consequently
the exterior constraint defect of $D(\tau)$, extended by zero to
the inside, is smooth and compactly supported: the beam vanishes
on an open collar of the sphere and outside its far shell.  Let
$\phi_\pm$ be these extended defects.

Fix a small preliminary loss $\epsilon>0$ and a base solve order
$s$ large enough for the later finite requirements.  The linear
beam estimates and finite-jet induction of hypersurface data imply
\begin{equation}\label{prep:seed-and-defects}
 \begin{split}
 \|D_\pm-(\delta,0)\|_{X^s}
        &\le C_s\rho_T^{-1}+A_Te^{\epsilon T},\\
 \|D_+-D_-\|_{X^s}&\le A_Te^{\epsilon T},\\
 \|\phi_\pm\|_{Y^s}&\le A_T^2e^{\epsilon T},\qquad
 \|\phi_+-\phi_-\|_{Y^s}\le A_T^3e^{\epsilon T}.
 \end{split}
\end{equation}
Here and below constants are absorbed by taking $T$ larger.  We
justify the defect estimates, including their derivative requirement.
For a metric on a spacelike hypersurface, Gauss--Codazzi expresses
the scalar and momentum constraints as the normal components of
its Einstein tensor.  Linearizing at the vacuum background, the
linear coefficient of the defect is therefore a linear combination
of the linear reduced-wave residual and one derivative of the linear
wave-gauge defect, with smooth background coefficients.  An $s$-jet
constraint bound thus requires gauge-defect jets through $s+1$,
in addition to the corresponding reduced-equation jets.  This
extra derivative is included before $\zeta$ is chosen.

More explicitly, Taylor expansion of this smooth finite-jet map
on the seed annulus gives
\[
 \phi(\tau)=\tau\ell_T+\tau^2Q_T+\tau^3R_T(\tau),
\]
where the linear term satisfies
$\|\ell_T\|_{Y^s}\le A_T\lambda_T^{-M_1}e^{o(T)}$, including
polynomial rescaling factors, and
\[
 \|Q_T\|_{Y^s}\le A_T^2\lambda_T^{P_s}e^{o(T)},\qquad
 \sup_{|\tau|\le1}\|R_T(\tau)\|_{Y^s}
           \le A_T^3\lambda_T^{P_s'}e^{o(T)}.
\]
The finite positive powers depend only on base jet orders.  The
quadratic term is identical at the two signs.  Taking $\zeta$
small for these powers and subsequently $M_1$ large enough to make
the linear term smaller than the cubic scale proves
Equation~\eqref{prep:seed-and-defects}.  This also accounts for the
nonlinear dependence of the second form on the induced lapse and
normal; it is not an assumption that the induced data depend
linearly on the sign.

Lemma~\ref{prep:constraint-pair} applies for large $T$.  Using its
same $\mathcal K$ for both signs yields
\begin{equation}\label{prep:correction-sizes}
 \|v_\pm\|_{X^s}\le C_s A_T^2e^{\epsilon T},\qquad
 \|v_+-v_-\|_{X^s}\le C_s A_T^3e^{2\epsilon T}.
\end{equation}
Undo Equation~\eqref{prep:rescaling} outside the ball, retain the
actual datum $d$ inside, and keep $d$ on the rest of the bridge.
Smoothness and vanishing of $v_\pm$ on the inner side make this a
smooth gluing.  The resulting $d_T^\pm$ satisfy the constraints
everywhere.  The physical interior is the original vacuum datum;
the artificial nonvacuum interior used by the solver has been
discarded and is never evolved.

For $s\ge m+3$, scaled Sobolev embedding bounds all required weighted
pointwise derivatives.  The cost of unscaling either correction
component is at most $C_m\rho_T$: for the metric the factors are
$\rho_T^{1+i}\rho_T^{-i}$, and for the second form they are
$\rho_T^{2+i}\rho_T^{-1-i}$.  Translation has bounded comparison
constants on the exterior region.  Thus
\begin{equation}\label{prep:physical-data-bound}
 p_m(d_T^\pm-d)
 \le \operatorname{poly}(T)\left(
       A_T\lambda_T^{m+5/2}e^{o(T)}
          +A_T^2e^{\epsilon T}\right).
\end{equation}
The extra seed derivative is the derivative entering the second
form.  Taking $\zeta$ small makes this tend to zero.  In particular
the metrics remain uniformly comparable to the complete original
metric, proving positivity and completeness.  The all-order symbol
bounds supplied by Lemma~\ref{prep:constraint-pair} give membership
in $\mathcal D$.  Finally
\[
 B\bigl(0,(1-2\beta-s_*)t_T-O(1)\bigr)\subset B_T;
\]
the protected physical region contains all fixed compact shadows
used in the finite bridge construction.  No outer cutoff is applied
to the correction; its allowed tail is part of $d_T^\pm$.

\subsection{Exact finite gauge jets}

Initialize the ordinary reduced equations on the partial exterior
slice $t_e=0$ using these physical data.  Its bounded bent portion
is in the development of unchanged data.  Prescribe the lapse and
shift to equal those of $g\pm h_T^B$ and use the corrected induced
metric and second form.  This fixes the full metric values and all
first spatial derivatives.  The formula for the second form fixes
the first temporal derivatives of the spatial metric entries.
The remaining first temporal derivatives are determined by
\begin{equation}\label{prep:wave-gauge}
 \Upsilon_\mu(\widetilde g)
 =-\widetilde g_{\mu\lambda}\widetilde g^{\alpha\beta}
       \bigl(\Gamma(\widetilde g)^\lambda_{\alpha\beta}
                         -\Gamma(g)^\lambda_{\alpha\beta}\bigr)=0.
\end{equation}
For clarity, in coordinates with initial surface $t_e=0$, put
$x_\nu=\partial_{t_e}\widetilde g_{0\nu}$.  After the spatial
temporal derivatives have been fixed, the lowered contracted
Christoffel expression is
\[
 \widetilde g^{\alpha\beta}
  \left(\partial_\alpha\widetilde g_{\beta\nu}
                   -\tfrac12\partial_\nu\widetilde g_{\alpha\beta}\right).
\]
Its coefficient on $x_0$ in the $\nu=0$ equation is
$\widetilde g^{00}/2$; the $x_i$ contributions there cancel.
The $\nu=i$ equations have coefficient $\widetilde g^{00}$ on
$x_i$.  Since the temporal slices are uniformly spacelike,
$\widetilde g^{00}$ is uniformly separated from zero.  This proves
the required smooth, uniformly invertible algebraic gauge solve.

Every higher initial temporal jet is recovered from the reduced
equation by solving for $\partial_{t_e}^2\widetilde g_{\mu\nu}$,
whose coefficient is again the nonzero $\widetilde g^{00}$, and
then differentiating.  Consequently, for each fixed finite order,
the complete initial spacetime jet is a smooth finite-jet function
of the corrected data and the prescribed lapse and shift.  Its
Taylor coefficients have the fixed-order bounds of the exterior
background.  Comparing it to the jets of $g+\tau h_T^B$, the linear
discrepancy is controlled by the linear gauge defect and reduced
residual and their finitely many derivatives.  The quadratic
discrepancy is even in $\tau$.  A correction contributes linearly
with a coefficient smooth in $\tau h_T^B$; the change of this
coefficient between signs costs another $A_Te^{\epsilon T}$.
Using Equation~\eqref{prep:correction-sizes}, the product of that
coefficient change and a quadratic correction is cubic.  Higher
Taylor terms have at least this order.  Thus, writing the exterior
tensor errors as
\[
 \mathfrak f_T^\pm=g_T^\pm-g\mp h_T^B,
\]
their initial full finite jets obey
\begin{equation}\label{prep:initial-gauge-errors}
 \|\mathfrak f_T^\pm\|_{\mathrm{init}}
      \le A_T^2e^{c\epsilon T},\qquad
 \|\mathfrak f_T^+-\mathfrak f_T^-\|_{\mathrm{init}}
      \le A_T^3e^{c\epsilon T}.
\end{equation}
The norms here can be any of the fixed ordinary jet or local slice
Sobolev norms on the finite computing domain; its polynomial volume
cost is included.  This argument applies to the recursively
recovered normal jets as well as the first jet.  On the protected
side the data, lapse, shift, and all gauge Cauchy jets agree exactly
with the background.

The solutions so initialized are vacuum wherever they exist.  For
the reduction
$\Ric(\widetilde g)+\nabla^{\widetilde g}_{(\mu}\Upsilon_{\nu)}=0$,
contracted Bianchi gives a homogeneous linear wave equation for
$\Upsilon$.  Its initial value is zero by
Equation~\eqref{prep:wave-gauge}.  The normal Einstein constraints
of the exact physical data then give its zero normal derivative
(the tangential derivatives of its zero initial value already
vanish).  Wave uniqueness propagates $\Upsilon=0$, hence
$\Ric(\widetilde g)=0$.  This uses the physical exterior and the
unchanged physical inner data only.

\subsection{Late support and the ordinary exterior rate}\label{prep:exterior-rate-section}

Work up to $t_e=t_T+C_2$ on the regular domain of
\cite[\CompanionRef{foundation}{pk2:exterior-comparison}]{Foundation}, with its deeper inner future exit and
the outer cutoff
\begin{equation}\label{prep:outer-cut}
 |y|<R_{\mathrm{out}}(t_e),\qquad
 R_{\mathrm{out}}(t_e)=C_{\mathrm{out}}(1+T)
                 +V_{\mathrm{out}}(t_T+C_2-t_e).
\end{equation}
The constants include the required spatial and time collars, and
$V_{\mathrm{out}}$ exceeds the uniform radial causal-speed bound.
Thus the outer face has strict future-outflow sign.  The domain is
only a computing restriction; the initial datum outside it is the
same full $d_T^\pm$, including its tail.

We first quantify support.  Suppose a causal path from the changed
initial support first reaches fixed large co-radius $R$ at time
$\tau\le t_T$.  Before that hit, its lab speed is at most
$1+\delta_R$, where $\delta_R\to0$ as $R\to\infty$, also in a
small comparison bootstrap.  Its initial point is outside $B_T$;
its hit point has the form $C_e(\tau)+O(R)$.  The triangle inequality
and the drift bound give
\begin{align}
 (1-2\beta)t_T
 &\le (1+\delta_R)\tau+s_*(t_T-\tau)+O(R),
                                   \label{prep:first-hit}\\
 t_T-\tau
 &\le \frac{2\beta+\delta_R}{1+\delta_R-s_*}\,t_T+O_R(1).
                                   \label{prep:late-hit}
\end{align}
For any prescribed $\theta>0$, first take $R$ large and then
$\beta$ small so that the last bound is at most $\theta T+O_R(1)$.
Only the inner exclusion radius occurs; no bound on the outer
support radius of the correction is used.  The beam and its
residuals have the same late compact-radius support by the return
ray bounds and tube margins in Proposition~\ref{bg:packet}.

These estimates also give exact early entry agreement.  At the
selected end, $I'\le C_I$ implies that
\[
 v<(1-\nu_0)T\quad\Longrightarrow\quad
 t_e\le t_T-(\nu_0/C_I)T+O(1).
\]
Take $\theta<\nu_0/(2C_I)$ and then $T$ large.  Finite propagation
and local reduced uniqueness from the exactly equal initial gauge
jets show that the exact change vanishes on this entry interval.
The entry is at bounded co-radius.  The bounded middle and
unselected pieces of its required interval lie in the unchanged
finite dependence region.  The deeper inner face is future outflow,
so it supplies no incoming path through the other end.  This proves
early agreement on the whole required entry interval.

For the size estimate, use lab components and write the ordinary
reduced equation for a change $H$ as
\begin{equation}\label{prep:exterior-equation}
 \mathcal R_e(H)_n=(\mathcal L_eH)_n
  +B^{\mu\nu}(H)\partial_\mu\partial_\nu H_n
  +\mathcal Q_n(H,\partial H),\qquad
 B(H)=(g+H)^{-1}-g^{-1}.
\end{equation}
A fixed background trace reversal can equivalently be included in
$\mathcal L_e$.  The smooth nonlinear terms start at degree two,
with total field derivative order at most two.  Indeed the
connection difference is a contraction of
$(g+H)^{-1}\nabla^gH$; the reduction leaves inverse-metric second
derivatives, first-derivative products, and algebraic background
curvature terms.  All required background coefficient jets cost
$e^{o(T)}$ on the regular domain.  At fixed orders the pure beam
source therefore obeys
\begin{equation}\label{prep:beam-nonlinear-residual}
 \|\mathcal R_e(\pm h_T^B)\|_{H^j_{\mathrm{mix}}}
 \le e^{o(T)}\left(A_T\lambda_T^{-M_1}
                           +A_T^2\lambda_T^{P_j}\right),
\end{equation}
including polynomial volume factors.  The positive $P_j$ are
fixed by the comparison orders, not the later formal accuracy.

Choose a high finite exterior energy order $k_e$.  Commute by mixed
lab derivatives and let $E_j^\pm(t_e)$, $0\le j\le k_e$, be the
square root of the sum of ordinary scalar-stress fluxes of the
order-$j$ components of $\mathfrak f_T^\pm$ in the actual metric
$g_T^\pm$.  Use the future slice normal as multiplier and a fixed
small positive mass $\mu$ in the energy.  Uniform regular geometry
makes these fluxes equivalent to
\[
 \sum_{|I|=j}\int\left(
   |\partial_{t_e}\partial^I\mathfrak f_T^\pm|^2
  +|\nabla_z\partial^I\mathfrak f_T^\pm|^2
  +\mu^2|\partial^I\mathfrak f_T^\pm|^2\right)dz.
\]
Bootstrap all these square-root energies by $e^{-\alpha T}$.
The beam base jets are exponentially small for sufficiently small
$\zeta$.  Sobolev embedding then preserves temporal
nondegeneracy and both outflow margins.  The slices admit uniformly
bounded local Sobolev extension and chart constants; the large
outer radius causes at most polynomial volume factors.

The exact perturbed principal operator is retained after every
commutation.  Every remaining nonlinear error term has another
packet or error factor, total field derivative order at most
$j+2$, and at most $j+1$ derivatives on an error factor.  Put the
highest error derivative in $L^2$.  Other error derivatives have
order at most $(j+2)/2$ and are controlled in supremum by the
higher fixed bootstrap energies.  Taking $k_e\ge16$ suffices for
these product placements at its top order.  A packet factor has
the fixed base bounds of Proposition~\ref{bg:packet}.  Thus these
terms contribute exponentially small coefficients times the
through-$j$ error energies.  There is no optical lapse loss in
this regular exterior calculation.

For the linear terms, the stress deformation and highest principal
commutators use first lab metric derivatives; first- and zeroth-order
tensor-wave terms use at most second derivatives.  These are bounded
globally and tend uniformly to zero at large co-radius by
Proposition~\ref{bg:exterior}.  Choose $\mu$ small first, then $R$
large: the same-level square-root rate is as small as prescribed
outside radius $R$.  Inside it, that rate is bounded by a constant
$C_{k_e,\mu}$ independent of the chosen $R$.  Further coefficient
hits multiply strictly lower error levels, with coefficients
$e^{o(T)}$.  All artificial boundary fluxes have the favorable sign.

Equation~\eqref{prep:late-hit} shows that the support sees the
possibly large rate only for $\theta T+O_R(1)$ time.  There is
therefore a nonnegative common majorant $r_e$ for both signs and
for the background linear problem such that, for every desired
$\epsilon_e>0$, the choices can give
\begin{equation}\label{prep:exterior-rate}
 \begin{split}
 (E_j^\pm)'&\le r_e E_j^\pm
       +b_T(t_e)\sum_{\ell<j}E_\ell^\pm+S_j,\\
 \mathcal R_e^{\mathrm{int}}(t)&:=\int_0^t r_e(s)\,ds,
 \qquad \mathcal R_e^{\mathrm{int}}(t_T+C_2)\le\epsilon_eT,
 \qquad \int_0^{t_T+C_2}b_T(s)\,ds=e^{o(T)}.
 \end{split}
\end{equation}
Here $S_j$ is bounded by Equation~\eqref{prep:beam-nonlinear-residual}.
Indeed the integral of the same-level rate is bounded by the small
far rate times $O(T)$ plus
$C_{k_e,\mu}(\theta T+O_R(1))$ and an exponentially small
bootstrap contribution.  Choosing these quantities in this order
proves the second line.  Enlarging $R$ does not enlarge the compact
rate constant.

To see explicitly the effect of higher coefficients, conjugate the
finite energy vector by $e^{-\mathcal R_e^{\mathrm{int}}(t)}$.
Its homogeneous couplings are strictly lower triangular.  The
time-ordered expansion terminates after $k_e$ strict drops; on
every subinterval $[s,t]$ its norm is bounded by
\[
 C_{k_e}\sum_{j=0}^{k_e}
       \left(1+\int_s^t b_T(\tau)\,d\tau\right)^j=e^{o(T)}.
\]
Thus the original propagator costs at most
$e^{\mathcal R_e^{\mathrm{int}}(t)-\mathcal R_e^{\mathrm{int}}(s)+o(T)}$.
Equation~\eqref{prep:initial-gauge-errors} and Duhamel's formula
give, with an arbitrarily small final loss $\epsilon_i>0$,
\begin{equation}\label{prep:individual-exterior}
 \max_{j\le k_e}E_j^\pm(t_e)\le A_T^2e^{\epsilon_iT}
 \qquad (0\le t_e\le t_T+C_2).
\end{equation}
This improves the bootstrap if $\epsilon_i<\alpha$.
Ordinary reduced-wave continuation therefore covers the finite
domain.  For each fixed $T$ its coefficients are nondegenerate;
inside the strict margins, one restarts the local Cauchy problem
and restricts by dependence.  The high base energy bounds prevent
finite-time breakdown.  Higher smooth norms propagate separately
by the differentiated linear estimates and require no simultaneous
smallness.  Exact gauge propagation, proved above, makes these
finite evolutions vacuum.

\subsection{The odd exterior estimate and the order of choices}

Subtract the two error equations, with the same background
$\mathcal L_e$ on the left.  The quadratic part of the nonlinear
expression in Equation~\eqref{prep:exterior-equation}, evaluated at
$h_T^B$ and $-h_T^B$, cancels exactly.  The remaining pure-beam
terms are cubic.  Every other nonlinear term contains an error
and at least one packet or further error factor.  At an order
$j$ at least a two-derivative and Sobolev buffer below $k_e$,
Equation~\eqref{prep:individual-exterior} consequently gives
\begin{equation}\label{prep:odd-exterior-source}
 \|\mathcal L_e(\mathfrak f_T^+-\mathfrak f_T^-)\|_{H^j_{\mathrm{mix}}}
 \le e^{c\epsilon_iT+o(T)}(A_T^3+A_T^4)
          +A_T\lambda_T^{-M_1}e^{o(T)}.
\end{equation}
In particular, terms such as a metric error times a second error
derivative use the two extra known individual derivatives.
No new quasilinear closure is needed for this lower-order equation.
Its support is inherited from the two actual solutions and the
beams, so the same late-support rate applies.  The initial odd jets
are cubic by Equation~\eqref{prep:initial-gauge-errors}.  Applying
the background version of Equation~\eqref{prep:exterior-rate}
therefore proves a cubic odd bound, with any prescribed small
loss, in all the required lower energy orders.

Transform the resulting full jets to the fixed background optical
coordinates on the entry collar and take their restriction to
$\mathcal S_T$.  The transition and normalized frame jets have
subexponential fixed-order bounds.  Slice Sobolev embedding with
two spare spatial derivatives gives the required uniform jets;
the spacetime collar supplies the normal jets.  The polynomial
entry volume is absorbed as well.  The full finite-jet matching
of $h_T^B$ to $h_T^p$ in Proposition~\ref{bg:packet} adds only the
linear inverse-accuracy error.  This proves
Equation~\eqref{prep:entry-estimates}.  The gauge condition is
tensorial with its fixed target $g$ and remains exact after this
coordinate change.  Exact early agreement was proved before the
energy estimate.  These are the actual entry jets used in the
interior Cauchy problem.

We finish by making the finite choices unambiguous.  Given $m,r$,
one may take
\[
 k_e\ge\max\{16,r+10\},\qquad
 s\ge\max\{m+3,k_e+8\},
\]
and then fix a single still larger finite beam base order covering
the hypersurface map, the constraint solve, the normal-jet
recursion, products, and traces.  Request the linear gauge-defect
bound with one derivative more than each constraint-defect order;
all these are finite requirements fixed at this stage.  If another
later finite residual order, such as the interior order $60$, is
needed, include it now.  Fix the small preliminary losses so that
their finite sum, including those in the odd comparison, is below
$\eta$.  Choose the mass, far-rate tolerance, radius threshold,
and shell tolerance so that both the exterior rate bound and the
given $\nu_0$ are respected.

Only now choose $\zeta>0$ sufficiently small for every positive
frequency power at these base orders and for
$(m+5/2)\zeta<\alpha$.  Next choose $M_1$ and the finite hierarchy,
beam, and matching degrees large enough that the linear residuals
are below the cubic scale at all these orders.  For example,
$M_1\zeta>4\kappa$ makes
$A_T\lambda_T^{-M_1}e^{o(T)}\le e^{-4\kappa T}$ for large $T$.
The linear construction permits this increase without increasing
the already fixed positive base powers.  Finally take $T$ large
so that the remaining fixed-order $e^{o(T)}$ factors fit the
reserved losses and all contractions and bootstraps apply.
This proves Proposition~\ref{prep:signed-entry}.  The all-order
smoothness of each resulting datum is separate from these finite
smallness choices; no one value of $\zeta$ is required to give
convergence at every derivative order simultaneously.

\begin{corollary}[Signed exterior preparation at a general amplitude]
\label{prep:amplitude-independent-exterior}
Fix any $\alpha>0$. Replace the amplitude in this section by
$A_T=e^{-\alpha T}$ and use Corollary~\ref{bg:amplitude-independent}.
For every fixed data order $m$, entry order $r$, loss $\eta>0$ and
late-support fraction $0<\nu_0<1$, the conclusions of
Proposition~\ref{prep:signed-entry} hold after choosing the frequency
exponent sufficiently small and the inverse accuracy sufficiently
large. This assertion ends at the entry cylinder.
\end{corollary}
\begin{proof}
The rescaled constraint solve and its odd difference estimates use
only $A_T\lambda_T^P\to0$ at finitely many fixed positive powers.
Choose the internal losses smaller than both the prescribed $\eta$
and a sufficiently small fixed fraction of $\alpha$. The exterior
bootstrap then improves because $A_T^2e^{\eta_{\rm int}T}=o(A_T)$.
Its common integrated rate in Equation~\eqref{prep:exterior-rate}
can be made arbitrarily small independently of $\alpha$, by the
same far-radius and first-arrival choices. The finite triangular
propagator preserves that rate and the odd cubic gain. Choose $\zeta$ so that every fixed positive frequency loss
$P\zeta$ fits its allocated part of the internal loss budget and is
smaller than $\alpha$. After fixing a positive margin $\delta$
within that budget, choose the finite inverse accuracy $M_1$ so that
\[
 \alpha+M_1\zeta>\max\{3\alpha,4\kappa\}+\delta.
\]
Then $A_T\lambda_T^{-M_1}e^{o(T)}$ is smaller than both the
cubic amplitude scale and $e^{-4\kappa T}$. These choices
complete the proof, with all-order smoothness obtained from the same
single base-order constraint solve. No interior comparison estimate
is used in this corollary.
\end{proof}
